% !TEX program = xelatex
% ============================================================
%  Arabic Exam Template — Nibras Center
%  Compile with: xelatex main.tex
%
%  NOTE ON DOCUMENT CLASS:
%    The `exam` document class is incompatible with the `bidi` package
%    (required for RTL Arabic via polyglossia). The `questions` environment
%    uses internal \if... conditionals that bidi breaks, causing
%    "Incomplete \iffalse" errors. This template therefore uses the
%    `article` class with manual exam-like commands that replicate the
%    key features: numbered questions, point values, choice lists,
%    answer boxes, a grade table, and an answer-reveal toggle.
%
%  ANSWER TOGGLE:
%    \printanswers    → shows correct answers in color
%    \noprintanswers  → hides answers (student version)
%  Switch the command below and recompile.
% ============================================================

\documentclass[11pt,a4paper]{article}

% ---- Geometry ----
\usepackage[a4paper,margin=2cm,top=2.2cm,bottom=2.2cm,headheight=15pt]{geometry}

% ---- Core packages (order matters: bidi before polyglossia) ----
\usepackage{amsmath}
\usepackage{graphicx}
\usepackage{enumitem}

% ---- RTL / Arabic language support ----
\usepackage{bidi}
\usepackage[no-math]{fontspec}
\usepackage[quiet]{polyglossia}

% ---- Fonts ----
\setmainfont[Scale=1]{NotoKufiArabic-Regular.ttf}
\newfontfamily\arabicfont[Script=Arabic,Scale=0.95]{NotoKufiArabic-Regular.ttf}
\newfontfamily\arabicfonttt[Script=Arabic,Scale=0.95]{NotoKufiArabic-Regular.ttf}
\newfontfamily\englishfont[Scale=0.95]{TeX Gyre Termes}

% ---- Languages ----
\setdefaultlanguage[calendar=gregorian,hijricorrection=1]{arabic}
\setotherlanguage[variant=british]{english}

% ---- Hyperref (load last, after polyglossia/bidi) ----
\usepackage[hidelinks]{hyperref}

% ---- Colors (table option for \rowcolor in grade table) ----
\usepackage[table]{xcolor}
\definecolor{answercolor}{HTML}{1B5E20}
\definecolor{pointcolor}{HTML}{1F3A5F}
\definecolor{gradehdr}{HTML}{1F3A5F}

% ---- Line spacing ----
\usepackage{setspace}
\setstretch{1.15}

% ---- Captions in Arabic ----
\addto\captionsarabic{%
  \renewcommand{\figurename}{الشكل}%
  \renewcommand{\tablename}{الجدول}%
}

% ============================================================
% ANSWER TOGGLE
% Change \noprintanswers to \printanswers to reveal answers.
% ============================================================
\newif\ifprintanswers
\newcommand{\printanswers}{\printanswerstrue}
\newcommand{\noprintanswers}{\printanswersfalse}
\noprintanswers          % student version (default)
\printanswers          % teacher version (uncomment to show answers)

% ============================================================
% Manual exam-class emulation
% ============================================================

% ---- Question counter ----
\newcounter{question}
\newcounter{choice}[question]

% \question[points] — starts a new numbered question
\newcommand{\question}[1][]{%
  \refstepcounter{question}%
  \par\vspace{0.8em}%
  \noindent\textbf{\thequestion.}%
  \if\relax\detokenize{#1}\relax\else\,\textcolor{pointcolor}{\textbf{(#1 درجة)}}\fi%
  \hspace{0.5em}%
}

% \part[points] — sub-part of a question (uses arabic letters)
\newcommand{\partitem}[1][]{%
  \stepcounter{choice}%
  \par\vspace{0.3em}%
  \noindent\hspace{1.5em}\textbf{\thequestion\,\alph{choice})}%
  \if\relax\detokenize{#1}\relax\else\,\textcolor{pointcolor}{\textbf{(#1 درجة)}}\fi%
  \hspace{0.4em}%
}

% \choice — one option in a multiple-choice list
\newcommand{\choice}{%
  \stepcounter{choice}%
  \par\noindent\hspace{2em}(\alph{choice})\hspace{0.4em}%
}

% \CorrectChoice — correct option; shown in color when \printanswers
\newcommand{\CorrectChoice}{%
  \stepcounter{choice}%
  \par\noindent\hspace{2em}(\alph{choice})\hspace{0.4em}%
  \ifprintanswers\color{answercolor}\bfseries\fi%
}

% \fillin[width] — fill-in-the-blank line
\newcommand{\fillin}[1][3cm]{%
  \ifprintanswers
    \colorbox{yellow!30}{\rule[-0.3em]{0pt}{1.6em}\hspace{#1}}%
  \else
    \underline{\hspace{#1}}%
  \fi%
}

% \fillwithlines{height} — lined answer space
\newcommand{\fillwithlines}[1]{%
  \par\vspace{0.3em}%
  \dimen0=#1\relax
  \loop
    \noindent\hrulefill\par\vspace{0.5em}%
    \advance\dimen0 by -1em\relax
  \ifdim\dimen0>1em\repeat
  \par
}

% \makeemptybox{height} — empty boxed answer space
\newcommand{\makeemptybox}[1]{%
  \par\vspace{0.3em}%
  \noindent\fbox{\rule{0pt}{#1}\rule{\dimexpr\linewidth-2\fboxsep-2\fboxrule\relax}{0pt}}%
  \par
}

% ---- Grade table (manual) ----
\newcommand{\gradetable}{%
  \par\vspace{0.3em}%
  \noindent
  \begin{tabular}{|c|c|c|c|c|c|c|}
    \hline
    \rowcolor{gradehdr}
    \textcolor{white}{\textbf{السؤال}} &
    \textcolor{white}{\textbf{1}} &
    \textcolor{white}{\textbf{2}} &
    \textcolor{white}{\textbf{3}} &
    \textcolor{white}{\textbf{4}} &
    \textcolor{white}{\textbf{5}} &
    \textcolor{white}{\textbf{6}} \\
    \hline
    \textbf{الدرجة} & & & & & & \\
    \hline
    \textbf{العلامة} & & & & & & \\
    \hline
  \end{tabular}%
  \par\vspace{0.3em}%
}

\begin{document}

% ============================================================
% TITLE BLOCK
% ============================================================
\begin{center}
{\Large\bfseries جامعة حلب --- جامعة حلب الخاصة}\\[0.3em]
{\large كلية الهندسة المعلوماتية}\\[0.3em]
{\large\bfseries امتحان مادة: \LR{Brain-Computer Interfaces}}\\[0.3em]
{\normalsize الفصل الأول \quad|\quad السنة الدراسية 2025/2026}\\[0.3em]
{\normalsize التاريخ: \today \quad|\quad المدة: 90 دقيقة}
\end{center}

\vspace{0.4em}
\hrule
\vspace{0.4em}

% ---- Student info fields ----
\noindent
\makebox[0.45\textwidth]{الاسم: \hrulefill} \hfill
\makebox[0.25\textwidth]{الرقم الجامعي: \hrulefill}

\vspace{0.3em}

\noindent
\makebox[0.30\textwidth]{الشعبة: \hrulefill} \hfill
\makebox[0.40\textwidth]{رقم الجلوس: \hrulefill}

\vspace{0.4em}
\hrule
\vspace{0.6em}

% ---- Grading grid ----
\gradetable

\vspace{0.6em}

% ============================================================
% QUESTIONS
% ============================================================

% ---- Q1: Multiple choice ----
\question[10]
ماذا يعني اختصار \LR{BCI}؟
\setcounter{choice}{0}
\CorrectChoice ترابط الدماغ–الحاسوب
\choice تخطيط كهربية الدماغ
\choice معالجة الإشارات الرقمية
\choice الذكاء الاصطناعي

% ---- Q2: True/False ----
\question[10]
حدد صحة العبارات التالية (صحيح / خطأ):
\setcounter{choice}{0}
\partitem[5] تُستخدم إشارات \LR{EEG} لقياس النشاط الكهربائي للدماغ.
\setcounter{choice}{0}
\CorrectChoice صحيح \choice خطأ

\setcounter{choice}{0}
\partitem[5] يُعتبر \LR{fMRI} أساساً غير جراحي لقراءة النشاط العصبي بدقة زمنية عالية.
\setcounter{choice}{0}
\choice صحيح \CorrectChoice خطأ

% ---- Q3: Fill in the blank ----
\question[10]
الوحدة الوظيفية الأساسية للجهاز العصبي هي \fillin[3cm]، وتُسمى بالعربية \fillin[4cm].

% TODO: ضع الإجابة الصحيحة هنا عند تفعيل \printanswers.

% ---- Q4: Short answer ----
\question[15]
اذكر ثلاثة من التحديات الرئيسية في تصميم أنظمة \LR{BCI}.
\fillwithlines{3cm}

% ---- Q5: Essay / long answer ----
\question[25]
اشرح بالتفصيل مراحل معالجة الإشارة في نظام \LR{BCI} نمطي، من الاكتساب حتى التصنيف.
% TODO: استبدل هذا السؤال بمحتوى الامتحان الفعلي.
\makeemptybox{5cm}

% ---- Q6: Short answer ----
\question[10]
قارن بين \LR{CSP} و \LR{LDA} من حيث الدور في خط معالجة \LR{BCI}.
\fillwithlines{3cm}

\vspace{1em}
\hrule
\vspace{0.3em}
\begin{center}
\textit{انتهى الامتحان --- بالتوفيق}
\end{center}

\end{document}
